% cf-macros.tex — 移植自 Colfax《Categorical Foundations for CuTe Layouts》官方源
% （arXiv:2601.05972，FinalVersion.tex L1-270 preamble）。
% 改动：①定理环境名中文化 ②xcolor 两次加载合并为一次 ③移除 lmodern/microtype
%（pdflatex 专属，xelatex 下不需要）④增加 longtable（术语表用）。
\usepackage{graphicx}
\usepackage{amsmath}
\usepackage{amsfonts}
\usepackage{amssymb}
\usepackage{amsthm}
\usepackage{tikz-cd}
\usetikzlibrary{matrix,fit,backgrounds}
\usepackage[table,dvipsnames]{xcolor}
\usepackage{mathrsfs}
\usepackage{mathdots}
\usepackage{mathtools}
\usepackage{comment}
\usepackage{array}
\usepackage{algorithm}
\usepackage{algpseudocode}  % breakable algorithmic environment
\usepackage{caption}
\usepackage[most]{tcolorbox}
\usepackage{longtable}
\usepackage{booktabs}

% Algorithm counter (for \ref)
\renewcommand{\thealgorithm}{\arabic{algorithm}}

% Breakable, framed algorithm box with continued title
\newtcolorbox[use counter=algorithm, number within=section]{BreakableAlgorithm}[2][]{%
  enhanced,
  breakable,
  colback=white,
  colframe=black,
  boxrule=0.8pt,
  arc=2pt,
  left=8pt,right=8pt,top=8pt,bottom=8pt,
  attach title to upper,
  title={算法~\thetcbcounter：#2},
  title after break={算法~\thetcbcounter（续）：#2},
  fonttitle=\bfseries,
  colbacktitle=white,
  coltitle=black,
  titlerule=0.8pt,
  borderline north={1pt}{0pt}{black},
  borderline south={1pt}{0pt}{black},
  #1
}

% Subtle line numbers & comment style
\algrenewcommand\alglinenumber[1]{\scriptsize\color{gray}{#1}}
\algrenewcommand{\algorithmiccomment}[1]{\hfill\(\triangleright\)\ #1}

% Horizontal separator you can place inside the algorithm
\newcommand{\algsep}{\Statex \rule{\linewidth}{0.6pt}}
\renewcommand{\baselinestretch}{1.10}
\usepackage[letterpaper,margin=1.2in]{geometry}  % 与原文版心一致（report 默认 letterpaper）
\setlength{\emergencystretch}{3em}  % CJK 行文断行点少，吸收长引用/行内数学小溢出
\usepackage[colorlinks=true, linkcolor=blue, urlcolor=blue, citecolor=blue]{hyperref}
\usepackage{cleveref}
\usepackage{listings}
\usepackage{enumitem}
\usepackage[backend=biber,style=numeric]{biblatex}
\addbibresource{references.bib}

\definecolor{colorx1}{rgb}{1.00, 0.00, 0.00} % Red
\definecolor{colorx2}{rgb}{1.00, 0.75, 0.00} % Orange-yellow
\definecolor{colorx3}{rgb}{0.50, 1.00, 0.00} % Yellow-green
\definecolor{colorx4}{rgb}{0.00, 1.00, 0.50} % Green-cyan
\definecolor{colorx5}{rgb}{0.00, 1.00, 1.00} % Cyan
\definecolor{colorx6}{rgb}{0.00, 0.50, 1.00} % Blue
\definecolor{colorx7}{rgb}{0.50, 0.00, 1.00} % Indigo
\definecolor{colorx8}{rgb}{1.00, 0.00, 1.00} % Magenta

\definecolor{codegray}{rgb}{0.5,0.5,0.5}
\definecolor{codebg}{rgb}{0.95,0.95,0.95}
\lstset{
    backgroundcolor=\color{codebg},
    basicstyle=\ttfamily\footnotesize,
    breaklines=true,
    frame=single,
    rulecolor=\color{codegray},
    keywordstyle=\color{blue},
    commentstyle=\color{gray},
    stringstyle=\color{red},
    showstringspaces=false,
}

\definecolor{keywordcolor}{rgb}{0.0, 0.0, 0.6}     % dark blue
\definecolor{identifiercolor}{rgb}{0.5, 0.0, 0.5}  % purple
\definecolor{stringcolor}{rgb}{0.58,0,0.82}        % magenta
\definecolor{bgcolor}{rgb}{0.95,0.95,0.95}         % light gray

% ===== COLORS =====
\definecolor{codebg}{rgb}{0.98, 0.98, 0.95}
\definecolor{codebase}{rgb}{0.25, 0.25, 0.28}
\definecolor{codecomment}{rgb}{0.45, 0.50, 0.53}
\definecolor{codekeyword}{rgb}{0.82, 0.37, 0.37}
\definecolor{codestring}{rgb}{0.37, 0.55, 0.29}
\definecolor{codefunc}{rgb}{0.48, 0.42, 0.93}
\definecolor{codetype}{rgb}{0.76, 0.47, 0.96}
\definecolor{codenumber}{rgb}{0.92, 0.47, 0.39}
\definecolor{codegray}{rgb}{0.60, 0.60, 0.60}
\definecolor{darkpastelpurple}{rgb}{0.59, 0.44, 0.84}
\definecolor{darkerpastelpurple}{rgb}{0.47, 0.34, 0.66}

% ===== CODE LISTINGS =====
\lstset{
    backgroundcolor=\color{codebg},
    basicstyle=\ttfamily\footnotesize\color{codebase},
    commentstyle=\itshape\color{codecomment},
    keywordstyle=\bfseries\color{keywordcolor},
    stringstyle=\color{codestring},
    identifierstyle=\color{darkerpastelpurple},
    numberstyle=\tiny\color{codegray},
    numbers=left,
    numbersep=6pt,
    frame=single,
    framesep=5pt,
    breaklines=true,
    showstringspaces=false,
    tabsize=4,
    morekeywords=[2]{cute, tract, uint32_t, int32_t, size_t, std, vector, string},
    keywordstyle=[2]\color{red},
    morekeywords=[3]{my_function, MyClass, CustomThing},
    keywordstyle=[3]\color{darkpastelpurple}\bfseries,
}

\lstdefinelanguage{cutedsl}{
    language=Python,
    morekeywords={},                     % 'cute' as a keyword
    morekeywords=[2]{},  % method names
    keywordstyle=\color{keywordcolor}\bfseries,
    morekeywords=[2]{cute, tract}
    keywordstyle=[2]\color{red},
    morestring=[b]',                        % single-quoted strings
    morestring=[b]",                        % double-quoted strings
    stringstyle=\color{stringcolor},
    showstringspaces=false,
    basicstyle=\ttfamily\footnotesize,
    backgroundcolor=\color{bgcolor},
    frame=single,
    breaklines=true,
    columns=flexible
}

% 定理环境（名称中文化，编号规则与原文一致：theorem 按 subsection 编号）
\newtheorem{theorem}{定理}[subsection]
\newtheorem{lemma}[theorem]{引理}
\newtheorem{corollary}[theorem]{推论}
\newtheorem{proposition}[theorem]{命题}
\newtheorem{mainthm}{主定理}
\renewcommand{\themainthm}{\Alph{mainthm}} % 主定理用大写字母 A, B, C...

\setcounter{secnumdepth}{3}

\theoremstyle{definition}
\newtheorem{definition}[theorem]{定义}
\newtheorem{example}[theorem]{例}
\newtheorem{construction}[theorem]{构造}
\newtheorem{observation}[theorem]{观察}
\newtheorem{notation}[theorem]{记号}
\newtheorem{warning}[theorem]{警告}
\newtheorem{localnotation}[theorem]{局部记号}

\theoremstyle{remark}
\newtheorem{remark}[theorem]{注}
\newtheorem{aside}[theorem]{旁注}

% amsthm proof 环境名 + cleveref 引用名中文化
\renewcommand{\proofname}{证明}
\crefname{theorem}{定理}{定理}
\crefname{lemma}{引理}{引理}
\crefname{corollary}{推论}{推论}
\crefname{proposition}{命题}{命题}
\crefname{mainthm}{主定理}{主定理}
\crefname{definition}{定义}{定义}
\crefname{example}{例}{例}
\crefname{construction}{构造}{构造}
\crefname{observation}{观察}{观察}
\crefname{notation}{记号}{记号}
\crefname{warning}{警告}{警告}
\crefname{localnotation}{局部记号}{局部记号}
\crefname{remark}{注}{注}
\crefname{aside}{旁注}{旁注}
\crefname{equation}{式}{式}
\crefname{figure}{图}{图}
\crefname{table}{表}{表}
\crefname{section}{节}{节}
\crefname{chapter}{章}{章}
\crefname{algorithm}{算法}{算法}

\newcommand{\m}{\langle m \rangle}
\newcommand{\n}{\langle n \rangle}
\newcommand{\p}{\langle p \rangle}
\newcommand{\q}{\langle q \rangle}

\newcommand{\fpset}[1]{\langle #1 \rangle_*}

\newcommand{\mathterm}[1]{\mathsf{#1}} % sans serif math terms

\newcommand{\sort}{\mathterm{sort}}
\newcommand{\coalesce}{\mathterm{coal}}
\newcommand{\coal}{\coalesce}
\renewcommand{\complement}{\mathterm{complement}}
\newcommand{\comp}{\mathterm{comp}}
\newcommand{\flatten}{\mathterm{flat}}
\newcommand{\cosize}{\mathterm{cosize}}
\newcommand{\size}{\mathterm{size}}
\newcommand{\len}{\mathterm{len}}
\newcommand{\rank}{\mathterm{rank}}
\newcommand{\depth}{\mathterm{depth}}
\newcommand{\mode}{\mathterm{mode}}
\newcommand{\entry}{\mathterm{entry}}
\newcommand{\red}{\mathterm{red}}
\newcommand{\sub}{\mathterm{sub}}
\newcommand{\shape}{\mathterm{shape}}
\newcommand{\stride}{\mathterm{stride}}
\newcommand{\colex}{\mathterm{colex}}
\newcommand{\squeeze}{\mathterm{squeeze}}
\newcommand{\Image}{\mathterm{Image}}
\newcommand{\codomain}{\mathterm{codomain}}
\newcommand{\domain}{\mathterm{domain}}
\newcommand{\complexity}{\mathterm{complexity}}
\newcommand{\id}{\mathterm{id}}
\newcommand{\incl}{\mathterm{incl}}
\newcommand{\Hom}{\mathterm{Hom}}
\renewcommand{\max}{\mathterm{max}}
\newcommand{\op}{\mathterm{op}}
\newcommand{\sk}{\mathterm{sk}}
\newcommand{\ob}{\mathterm{ob}}
\newcommand{\filter}{\mathterm{filter}}
\newcommand{\full}{\mathterm{full}}
\newcommand{\profile}{\mathterm{prof}}
\newcommand{\mor}{\mathterm{mor}}
\newcommand{\row}{\mathterm{row}}
\newcommand{\col}{\mathterm{col}}
\newcommand{\tiled}{\mathterm{tiled}}
\newcommand{\ind}{\mathterm{ind}}

\newcommand{\bit}[1]{\textbf{\textit{#1}}}
\newcommand{\catstyle}[1]{{\boldsymbol{\mathsf{#1}}}}
\newcommand{\msf}[1]{\mathsf{#1}}

\newcommand{\Set}{\catstyle{Set}}
\newcommand{\Fin}{\catstyle{Fin}}
\newcommand{\FinSet}{\catstyle{FinSet}}
\newcommand{\Tuple}{\catstyle{Tuple}}
\newcommand{\Cat}{\catstyle{Cat}}
\newcommand{\cC}{\catstyle{C}}
\newcommand{\cD}{\catstyle{D}}
\newcommand{\cE}{\catstyle{E}}
\newcommand{\cS}{\catstyle{S}}
\newcommand{\Poset}{\catstyle{Poset}}
\newcommand{\NN}{\mathbb{N}}
\newcommand{\adj}{\dashv}

\definecolor{amethyst}{rgb}{0.6, 0.4, 0.8}
\newcommand{\rs}[1]{{\color{amethyst}\textbf{RS: }#1}}
